\section{SDE inversas y proceso generativo}

\subsection{Teorema de Anderson: SDE inversa}

Dado la SDE hacia adelante

$$
\mathrm{d}\mathbf{x} = f(\mathbf{x}, t)\,\mathrm{d}t + g(t)\,\mathrm{d}\mathbf{w}
$$

Anderson (1982) demostró que la SDE de tiempo inverso correspondiente es:

$$
\mathrm{d}\mathbf{x} = \left[f(\mathbf{x}, t) - g(t)^2 \nabla_{\mathbf{x}} \log p_t(\mathbf{x})\right]\mathrm{d}t + g(t)\,\mathrm{d}\bar{\mathbf{w}}
$$

donde $\bar{\mathbf{w}}$ es un proceso de Wiener en tiempo inverso, y $\mathrm{d}t$ representa aquí una \textbf{incremento negativo} de tiempo (el tiempo avanza desde $T$ hacia 0).

\begin{importantbox}{Estructura de la SDE inversa}
La SDE inversa comparte una estructura sorprendentemente similar a la SDE hacia adelante:
\begin{itemize}
    \item El coeficiente de difusión $g(t)$ es \textbf{idéntico};
    \item El término de arrastre resta \textbf{adicionalmente} $g(t)^2 \nabla_{\mathbf{x}} \log p_t(\mathbf{x})$ sobre la base del original $f(\mathbf{x}, t)$.
\end{itemize}
Lo único que se requiere conocer de manera adicional es la \textbf{función score} $\nabla_{\mathbf{x}} \log p_t(\mathbf{x})$, que es precisamente el objetivo aprendido por las redes neuronales en los modelos de difusión.
\end{importantbox}

\subsection{El papel de la función Score}

Repasando el enfoque utilizado en DDPM discreto:
\begin{enumerate}
    \item Definir el proceso hacia adelante (añadir ruido);
    \item Entrenar una red de predicción de ruido $\boldsymbol{\epsilon}_\theta(\mathbf{x}_t, t)$;
    \item Obtener la función score mediante la relación $\nabla_{\mathbf{x}} \log p_t(\mathbf{x}_t) = -\frac{\boldsymbol{\epsilon}}{\sigma_t}$;
    \item Ejecutar el proceso inverso de desruido utilizando la función score.
\end{enumerate}

En un marco de tiempo continuo, la lógica es totalmente consistente:
\begin{enumerate}
    \item Definir la SDE hacia adelante (arrastre + difusión);
    \item Entrenar una red score $\mathbf{s}_\theta(\mathbf{x}, t) \approx \nabla_{\mathbf{x}} \log p_t(\mathbf{x})$;
    \item Sustituir $\mathbf{s}_\theta$ en la SDE inversa para ejecutar la generación.
\end{enumerate}

\begin{knowledgebox}{Equivalencia entre score y predicción de ruido}
Para una distribución de salto $q(\mathbf{x}_t | \mathbf{x}_0) = \mathcal{N}(\alpha_t \mathbf{x}_0, \sigma_t^2 \mathbf{I})$, la función score puede expresarse como:
$$
\nabla_{\mathbf{x}_t} \log p_t(\mathbf{x}_t | \mathbf{x}_0) = -\frac{\mathbf{x}_t - \alpha_t \mathbf{x}_0}{\sigma_t^2} = -\frac{\boldsymbol{\epsilon}}{\sigma_t}
$$
donde $\boldsymbol{\epsilon}$ es el ruido gaussiano estándar muestreado al añadir ruido. Por lo tanto, predecir la función score es equivalente a predecir el ruido (la diferencia radica en dividir por $\sigma_t$). Esto explica por qué el objetivo de entrenamiento de "predicción de ruido" en DDPM y el objetivo de Score Matching son esencialmente lo mismo.
\end{knowledgebox}

\subsection{Implementación discreta del proceso inverso}

Aunque la SDE inversa es continua, se requiere discretizarla para el cálculo práctico. Dada una SDE inversa, dividimos el tiempo $[0, T]$ en $N$ pasos; en cada paso:

$$
\mathbf{x}_{t-\Delta t} \approx \mathbf{x}_t + \left[f(\mathbf{x}_t, t) - g(t)^2 \mathbf{s}_\theta(\mathbf{x}_t, t)\right](-\Delta t) + g(t)\sqrt{\Delta t}\;\boldsymbol{\epsilon}
$$

Esto es totalmente consistente con el proceso inverso discreto de DDPM: en cada paso se corrige la media utilizando la función score (o predicción de ruido) y luego se añade un ruido aleatorio apropiado.

\begin{warningbox}{Número de pasos de discretización y calidad de generación}
La discretización de una SDE es esencialmente el método de Euler-Maruyama. Cuantos más pasos, más precisa será la aproximación y mejor será la calidad de generación; sin embargo, también aumentará el costo computacional. DDPM requiere aproximadamente 1000 pasos para obtener una buena calidad, mientras que DDIM (la discretización correspondiente a una ODE determinista) solo necesita aproximadamente 50 pasos. En conferencias posteriores se presentarán solvers de SDE/ODE más eficientes (como DPM-Solver), lo cual permitirá reducir aún más el número de pasos a un rango de 10--15.
\end{warningbox}

\subsection{Resumen de la sección}

El teorema de Anderson garantiza que: dada cualquier SDE hacia adelante, la SDE inversa se determina automáticamente; lo único que requiere aprendizaje es la función score. Una vez discretizada la SDE inversa, recuperamos el algoritmo de muestreo de los modelos de difusión clásicos.

